\documentclass[preview]{standalone}

\usepackage{amsmath}
\usepackage{amssymb}
\usepackage{stellar}
\usepackage{definitions}

\begin{document}

\id{psicologia-psicofisica}
\genpage

\section{Psicofisica}

\begin{snippetdefinition}{psicofisica-definition}{Psicofisica}
    La \emph{psicofisica} è lo studio della relazione tra gli stimoli fisici e
    il comportamento o le esperienze mentali evocate da tali stimoli.
\end{snippetdefinition}

\begin{snippet}{fechner-psicofisica}
    Gustav Fechner coniò il termine \emph{psicofisica}. Il problema centrale
    della psicofisica consiste nel misurare l'intensità delle esperienze
    sensoriali e nel collegarla all'intensità fisica degli stimoli.
\end{snippet}

\begin{snippetdefinition}{soglia-assoluta-definition}{Soglia assoluta}
    La \emph{soglia assoluta} della stimolazione è la quantità minima di energia
    fisica necessaria per produrre un'esperienza sensoriale rilevata nel
    \(50\%\) dei casi.
\end{snippetdefinition}

\begin{snippetexample}{soglia-assoluta-example}{Soglia assoluta}
    Per misurare una soglia assoluta uditiva, i ricercatori possono chiedere ai
    partecipanti di ascoltare suoni molto deboli in una stanza insonorizzata.
    In una serie di prove vengono presentati stimoli di diversa intensità e, in
    ogni prova, il partecipante indica se ha rilevato lo stimolo.
\end{snippetexample}

\includesnpt[width=62\%,src=/snippet/static-psychology/psychophysical-detection-function.jpg]{centered-img}

\begin{snippetdefinition}{abituazione-sensoriale-definition}{Abituazione sensoriale}
    L'\emph{abituazione sensoriale} è la diminuzione della responsività di un
    sistema sensoriale sottoposto a una stimolazione prolungata.
\end{snippetdefinition}

\begin{snippetdefinition}{distorsioni-giudizio-definition}{Distorsioni di giudizio}
    Le \emph{distorsioni di giudizio}, o \emph{bias}, sono tendenze sistematiche
    dell'osservatore a rispondere in un modo particolare per ragioni non
    dipendenti dalle caratteristiche sensoriali dello stimolo.
\end{snippetdefinition}

\begin{snippetdefinition}{teoria-detezione-segnale-definition}{Teoria della detezione del segnale}
    La \emph{teoria della detezione del segnale}, o \emph{TDS}, è un approccio
    psicofisico che distingue il ruolo della sensibilità allo stimolo dal ruolo
    del criterio individuale di risposta nel giudicare la presenza o l'assenza
    di uno stimolo.
\end{snippetdefinition}

\begin{snippet}{risposte-detezione-segnale}
    Nella teoria della detezione del segnale, alcune prove contengono uno
    stimolo debole e altre non contengono alcuno stimolo. In ogni prova,
    l'osservatore risponde se ritiene che il segnale sia presente. Le risposte
    possibili sono:
    \begin{itemize}
        \item \emph{hit}: il soggetto indica che il segnale è presente quando
        esso è effettivamente comparso;
        \item \emph{omissione}: il soggetto indica che il segnale non è
        presente quando esso è comparso;
        \item \emph{falso allarme}: il soggetto indica che il segnale è
        presente quando esso non è comparso;
        \item \emph{rifiuto corretto}: il soggetto indica che il segnale è
        assente quando esso non è comparso.
    \end{itemize}
\end{snippet}

\includesnpt[width=62\%,src=/snippet/static-psychology/signal-detection-theory-outcomes.jpg]{centered-img}

\begin{snippetdefinition}{soglia-differenziale-definition}{Soglia differenziale}
    La \emph{soglia differenziale} è la più piccola differenza fisica tra due
    stimoli che può essere riconosciuta come differenza minima rilevabile nel
    \(50\%\) dei casi.
\end{snippetdefinition}

\begin{snippetdefinition}{differenza-minima-rilevabile-definition}{Differenza minima rilevabile}
    La \emph{differenza minima rilevabile} è un'unità quantitativa che misura la
    grandezza della differenza psicologica tra due sensazioni.
\end{snippetdefinition}

\begin{snippetexample}{soglia-differenziale-example}{Soglia differenziale}
    Una compagnia di bibite vuole produrre una cola percepita come più dolce
    rispetto a quelle già presenti sul mercato, riducendo al minimo la quantità
    di zucchero utilizzata. In questo caso occorre misurare la soglia
    differenziale, cioè la minima differenza nella quantità di zucchero che può
    essere percepita.
\end{snippetexample}

\begin{snippetdefinition}{legge-weber-definition}{Legge di Weber}
    La \emph{legge di Weber} afferma che la differenza minima rilevabile tra
    due stimoli è una frazione costante dell'intensità dello stimolo standard:
    \[
        \frac{\Delta I}{I} = K
    \]
    dove \(K\) è la costante di Weber, \(I\) è l'intensità dello stimolo
    standard e \(\Delta I\) è l'incremento che produce una differenza minima
    rilevabile.
\end{snippetdefinition}

\begin{snippetexample}{legge-weber-example}{Legge di Weber}
    In un esperimento sulla lunghezza, si può chiedere ai partecipanti di
    indicare se due linee hanno la stessa lunghezza. Weber osservò che ogni
    modalità sensoriale presenta un valore specifico di \(K\); per la
    percezione della lunghezza, ad esempio, il valore può essere circa
    \(0{,}1\).
\end{snippetexample}

\includesnpt[width=62\%,src=/snippet/static-psychology/webers-law-perceptual-comparison.jpg]{centered-img}

\begin{snippetdefinition}{legge-fechner-definition}{Legge di Fechner}
    La \emph{legge di Fechner} afferma che l'intensità della sensazione è
    direttamente proporzionale al logaritmo dell'intensità dello stimolo:
    \[
        S = K \log I + C
    \]
    dove \(S\) è l'intensità della sensazione, \(K\) è la costante di Weber,
    \(I\) è l'intensità dello stimolo e \(C\) è una costante di integrazione.
\end{snippetdefinition}

\begin{snippet}{fechner-relazione-logaritmica}
    Secondo Fechner, se lo stimolo aumenta secondo una progressione geometrica,
    la sensazione corrispondente aumenta secondo una progressione aritmetica.
\end{snippet}

\includesnpt[width=62\%,src=/snippet/static-psychology/fechner-logarithmic-law.jpg]{centered-img}

\end{document}
